%% ============================================================
%%  latam-es/experience-es.tex  —  Experiencia (LATAM · conciliado)
%%
%%  Historial conciliado contra el Certificado Único Laboral
%%  (MTPE Perú, emitido el 2026-10-26).
%%
%%  Regla de conciliación: cuando dos registros de planilla se solapan,
%%  gana el contrato más extenso; el más corto se recorta al espacio libre
%%  y se descarta si queda por debajo de 25 días. Cero solapes por
%%  construcción.
%%
%%  Las fechas vienen del CUL salvo en los tres empleadores extranjeros
%%  (ExSquared, DaCodes, The Bridge Social), que no pueden figurar en
%%  planilla peruana.
%%
%%  Descartados por la regla: Kruger Corp / BCP, NTT Data / Telefónica,
%%  Tata / SUNAT, Caleidos / Apple, Quality & Development / AUNA,
%%  Neoris / Cencosud, Q System / Inkafarma, AVLA y el 2º contrato de
%%  Manpower.
%%
%%  Descartada por el candidato (2026-10-26): Alternativa Tecnológica / BCP.
%%  La ventana mar 2015 – jul 2017 donde estaba la cubre el bloque freelance.
%%
%%  Fusionadas (2026-10-26): WProducciones y LimaW son la misma empresa, así
%%  que sus dos registros de planilla entran como una sola entrada de 10 meses.
%%
%%  13 roles — la carrera abre en mayo 2014, no en enero 2014.
%%
%%  Apple Inc. queda FUERA de este variante a propósito: va solo en el
%%  variante USA (usa/experience.tex).
%% ============================================================

\section{Experiencia Profesional}

%% ── 1. IDM Technology / Scotiabank  (mar 2026 – presente) ───
\cvEntryClient%
  {Desarrollador Fullstack Senior}%
  {IDM Technology --- Lima, Perú}%
  {Scotiabank Perú}%
  {Plataforma de Pagos QR Dinámico}%
  {Marzo 2026 -- Presente}

\begin{cvItems}
  \cvItem{Construyó microservicios reactivos cloud-native (\textbf{Java 25, Spring Boot 4,
    Spring WebFlux}) para generación de QR dinámico y estático, procesamiento de pagos y
    persistencia transaccional vía \textbf{APIs REST}.}
  \cvItem{Diseñó \textbf{servicios de integración} conectando \textbf{Google Cloud Platform} y
    \textbf{Apigee} con sistemas mainframe IBM AS/400 mediante DataQueues.}
  \cvItem{Implementó autenticación de grado bancario con \textbf{Spring Security} y optimizó las
    capas de persistencia en \textbf{SQL Server} para cargas transaccionales de alto volumen.}
  \cvItem{Desarrolló herramientas internas para la operación del banco con \textbf{React},
    \textbf{Electron}, \textbf{Node.js} y \textbf{Python}.}
  \cvItem{Mantuvo pipelines \textbf{CI/CD en Jenkins} y quality gates de \textbf{SonarCloud};
    remedió vulnerabilidades detectadas en pruebas de penetración en equipos Agile Scrum.}
  \cvItem{Aplicó \textbf{desarrollo asistido por IA} (Claude Code, GitHub Copilot) para acelerar
    scaffolding, revisión de código y entrega.}
\end{cvItems}

{\footnotesize\color{cvMuted}\textit{Stack:}~%
Java 25, Spring Boot 4, Spring WebFlux, APIs REST, SQL Server, IBM AS/400, DataQueues,
Spring Security, GCP, Apigee, React, Electron, Node.js, Python, Jenkins, SonarCloud, Bitbucket.}

\vspace{2pt}

%% ── 2. DaCodes / HealthCare — AgentCloud  (ene 2025 – feb 2026)
\cvEntryClient%
  {Ingeniero de Software Líder}%
  {DaCodes --- Remoto, Estados Unidos}%
  {HealthCare}%
  {AgentCloud}%
  {Enero 2025 -- Febrero 2026}

\begin{cvItems}
  \cvItem{Lideró la arquitectura y el desarrollo full-stack de \textbf{AgentCloud}, una
    \textbf{plataforma SaaS} cloud-native para una aseguradora de salud estadounidense ---
    gestionando asegurados, proveedores médicos y redes de clínicas de extremo a extremo.}
  \cvItem{Construyó microservicios de \textbf{afiliación de asegurados}, gestión de pólizas,
    cobertura de beneficios y flujos de autorización de atención sobre \textbf{AWS}
    (Lambda, EKS, SQS) con \textbf{Java / Spring Boot} y \textbf{Scala}; contenerizados con Docker
    y orquestados en Kubernetes.}
  \cvItem{Desarrolló el portal administrativo en \textbf{React 18} para que el personal de la
    aseguradora gestione asegurados, clínicas, planes de cobertura, autorizaciones y tableros de
    reportes en tiempo real.}
  \cvItem{Construyó la aplicación móvil en \textbf{React Native} para que los asegurados consulten
    su cobertura, ubiquen clínicas de la red y soliciten autorizaciones desde su teléfono.}
  \cvItem{Implementó el backend \textbf{orientado a eventos} con \textbf{Apache Kafka} para
    procesamiento asíncrono de autorizaciones, notificaciones de siniestros y disparadores de
    facturación de proveedores.}
  \cvItem{Diseñó y optimizó interacciones con \textbf{Oracle 19c}, incluyendo vistas complejas y
    procedimientos almacenados ajustados al rendimiento.}
  \cvItem{Creó \textbf{microservicios serverless en Python} sobre AWS Lambda para procesamiento
    documental y tareas programadas.}
  \cvItem{Definió estándares de ingeniería, lideró revisiones de código y guió al equipo
    distribuido en ceremonias Agile; CI/CD con GitHub Actions. Aplicó \textbf{desarrollo asistido
    por IA} (Claude Code, GitHub Copilot, Cursor).}
\end{cvItems}

{\footnotesize\color{cvMuted}\textit{Stack:}~%
React 18, React Native, Scala, Java, Spring Boot, Python, Oracle 19c, Apache Kafka,
AWS (S3, Lambda, RDS, EKS, SQS, CloudWatch), Docker, Kubernetes, GitHub Actions.}

\vspace{2pt}

%% ── 3. ExSquared / SambaSafety — Qorta  (mar 2023 – dic 2024) ─
\cvEntryClient%
  {Ingeniero Full Stack Senior}%
  {ExSquared (antes Outcoding) --- Remoto, Estados Unidos}%
  {SambaSafety}%
  {Qorta --- Driver Manager}%
  {Marzo 2023 -- Diciembre 2024}

\begin{cvItems}
  \cvItem{Contribuyó a \textbf{Qorta}, la plataforma de monitoreo y analítica de riesgo de
    conductores y vehículos de SambaSafety, usada por operadores de flotas en Norteamérica.}
  \cvItem{Desarrolló y migró módulos frontend a \textbf{micro frontends} con \textbf{Angular 15}
    y \textbf{Tailwind CSS}.}
  \cvItem{Construyó microservicios con \textbf{Java 17} y \textbf{Spring Boot} siguiendo
    \textbf{arquitectura hexagonal}.}
  \cvItem{Sustituyó los listeners heredados de \textbf{RabbitMQ} por \textbf{Spring Cloud Stream}
    en los microservicios existentes.}
  \cvItem{Implementó flujos de procesamiento de datos con \textbf{Apache Kafka} y
    \textbf{Redis}.}
  \cvItem{Construyó \textbf{APIs REST} e integró proveedores de registros vehiculares (MVR) de
    múltiples jurisdicciones de EE.\,UU.; optimizó consultas \textbf{PostgreSQL} sobre conjuntos
    masivos de datos de conductores.}
  \cvItem{Mantuvo \textbf{95\%+ de cobertura} en frontend y backend con \textbf{JUnit},
    \textbf{Mockito} y \textbf{Jest} para cumplir los quality gates del CI/CD.}
\end{cvItems}

{\footnotesize\color{cvMuted}\textit{Stack:}~%
Java 17, Spring Boot, Spring Cloud, Spring Data, Angular 15, TypeScript, Tailwind, Sass, Jest,
PostgreSQL, Apache Kafka, Redis, JWT, AWS (EKS, SES, SNS, S3), Git, Bitbucket.}

\vspace{2pt}

%% ── 4. Globant / Banco GNB Paraguay  (jun 2022 – feb 2023) ───
\cvEntryClient%
  {Ingeniero Full Stack Senior}%
  {Globant --- Remoto, Paraguay}%
  {Banco GNB Paraguay}%
  {Banca Personal y Corporativa}%
  {Junio 2022 -- Febrero 2023}

\begin{cvItems}
  \cvItem{Entregó funcionalidades de banca digital como ingeniero embebido de Globant en el equipo
    de producto del Banco GNB Paraguay, cubriendo gestión de cuentas, pago de servicios y la
    operación completa de banca personal y empresarial.}
  \cvItem{Desarrolló y mantuvo microservicios con \textbf{Java 11 / Java 8} y \textbf{Spring
    Boot}, cumpliendo requisitos de escalabilidad, seguridad y rendimiento.}
  \cvItem{Desarrolló la aplicación móvil de banca en \textbf{React Native} para gestión de
    cuentas, historial de transacciones y notificaciones push en tiempo real --- publicada en
    App Store y Google Play.}
  \cvItem{Construyó nuevos componentes en \textbf{Angular 12} para los flujos de banca por
    internet del cliente individual y empresarial, con estilos en CSS3 y pruebas en \textbf{Jest}.}
  \cvItem{Implementó lógica de negocio y persistencia con \textbf{JPA / Hibernate}; integró
    servicios REST externos mediante \textbf{Feign}.}
  \cvItem{Lideró la migración de microservicios heredados y mantuvo \textbf{85\%+ de cobertura}
    con \textbf{JUnit} y \textbf{Mockito}, incluyendo excepciones controladas, para cumplir el
    gate de CI/CD.}
\end{cvItems}

{\footnotesize\color{cvMuted}\textit{Stack:}~%
Java 11, Java 8, Spring Boot, Spring Cloud, Spring Data, Angular 12, React Native, JavaScript,
CSS3, Jest, SQL Server, JPA/Hibernate, Feign, JWT, AWS (EKS), Git, GitLab.}

\vspace{2pt}

%% ── 5. The Bridge Social / NCB — Lynk  (ene 2021 – may 2022) ─
\cvEntryClient%
  {Ingeniero Full Stack}%
  {The Bridge Social --- Remoto, Jamaica}%
  {National Commercial Bank}%
  {Lynk}%
  {Enero 2021 -- Mayo 2022}

\begin{cvItems}
  \cvItem{Construyó \textbf{Lynk}, plataforma nacional de banca digital y monedero móvil del
    National Commercial Bank de Jamaica, integrando sistemas core bancarios internos y terceros.}
  \cvItem{Desarrolló aproximadamente \textbf{50 microservicios} con \textbf{Java 11} y
    \textbf{Spring Boot}, aplicando \textbf{arquitectura hexagonal} por la cantidad de fuentes de
    información involucradas.}
  \cvItem{Implementó \textbf{arquitectura orientada a eventos} con \textbf{Apache Kafka}; integró
    servicios mediante \textbf{Feign} y \textbf{RestTemplate}.}
  \cvItem{Integró proveedores externos, entre ellos \textbf{Twilio} para flujos de mensajería y
    verificación.}
  \cvItem{Construyó \textbf{microservicios serverless} con \textbf{Azure Functions} y
    \textbf{Node.js / NestJS}.}
  \cvItem{Desarrolló el frontend en \textbf{React} con \textbf{Tailwind CSS} y las pantallas de
    monedero móvil en \textbf{React Native}; pruebas unitarias con \textbf{Jest}.}
  \cvItem{Entregó sobre \textbf{Azure} (AKS, App Insights, Event Hubs, Azure SQL) con
    \textbf{PostgreSQL} como motor relacional; aplicó patrones builder, singleton y otros.}
  \cvItem{Gestionó métricas de calidad con \textbf{SonarCloud} y el ciclo completo del proyecto en
    \textbf{Azure DevOps} con \textbf{Azure Repos}.}
\end{cvItems}

{\footnotesize\color{cvMuted}\textit{Stack:}~%
Java 11, Spring Boot, Spring Cloud, Spring Data, Node.js, NestJS, Apache Kafka, React,
React Native, Tailwind, Jest, PostgreSQL, JWT, Azure (AKS, Functions, Event Hubs, App Insights,
Azure SQL), SonarCloud, Azure DevOps.}

\vspace{2pt}

%% ── 6. Indra Perú / RIMAC — Póliza Electrónica  (abr – dic 2020)
%%     Contenido aportado por el candidato el 2026-10-26.
%%     TODO menor: confirmar el título exacto del rol.
\cvEntryClient%
  {Ingeniero Full Stack Senior}%
  {Indra Perú --- Lima, Perú}%
  {RIMAC Seguros}%
  {Póliza Electrónica --- Nuevos Productos}%
  {Abril 2020 -- Diciembre 2020}

\begin{cvItems}
  \cvItem{Construyó nuevos productos de seguros para \textbf{RIMAC Seguros}, entre ellos
    \textbf{Póliza Electrónica}, digitalizando la emisión de pólizas de extremo a extremo.}
  \cvItem{Asumió la capa de datos completa en \textbf{Oracle} con \textbf{PL/SQL}: procedimientos
    almacenados, packages, vistas y el modelo que soporta la emisión y la configuración de
    productos.}
  \cvItem{Desarrolló \textbf{servicios serverless} en \textbf{AWS Lambda} con \textbf{Node.js},
    usando \textbf{S3} para el almacenamiento documental junto a otros servicios de AWS.}
  \cvItem{Construyó los frontends de producto en \textbf{React}, cubriendo los flujos de
    configuración y emisión para usuarios internos y para corredores.}
\end{cvItems}

{\footnotesize\color{cvMuted}\textit{Stack:}~%
Oracle, PL/SQL, Node.js, React, AWS (Lambda, S3), APIs REST.}

\vspace{2pt}

%% ── 7. Michael Page / Intercorp Retail  (ene – mar 2020) ─────
%%     Recortado. El CUL registra 26/12/2019 -- 31/03/2020; Zoluxiones
%%     ocupa diciembre de 2019.
\cvEntryClient%
  {Ingeniero Full Stack}%
  {Michael Page International --- Lima, Perú}%
  {Intercorp Retail}%
  {Tracking de Incidentes}%
  {Enero 2020 -- Marzo 2020}

\begin{cvItems}
  \cvItem{Diseñó la arquitectura general de una plataforma que mapea las quejas de clientes de
    todas las empresas del grupo Intercorp --- almacenando pruebas y notificando a los
    involucrados, integrada con el stack tecnológico distinto de cada subsidiaria.}
  \cvItem{Desarrolló microservicios en \textbf{Java 8} con \textbf{Spring Boot} y en
    \textbf{Node.js} con \textbf{NestJS}; normalizó el esquema \textbf{PostgreSQL}.}
  \cvItem{Integró \textbf{Google Cloud Storage} para la carga de archivos y \textbf{AWS SES} para
    el envío de correos a clientes y gestores.}
  \cvItem{Desarrolló microservicios de autenticación con \textbf{JWT} y \textbf{Spring Security}.}
  \cvItem{Definió la estructura arquitectónica de la aplicación \textbf{Angular 8} y colaboró con
    diseñadores UX/UI para traducir los requisitos de diseño en componentes y funcionalidades,
    con \textbf{Angular Material} y \textbf{Bootstrap}.}
  \cvItem{Desplegó servicios contenerizados en \textbf{Google Cloud Platform}; evaluó cobertura
    con \textbf{SonarQube} y gestionó la entrega en \textbf{Azure DevOps} bajo Scrum.}
\end{cvItems}

{\footnotesize\color{cvMuted}\textit{Stack:}~%
Java 8, Spring Boot, Node.js, NestJS, Angular 8, Angular Material, Bootstrap, PostgreSQL, JWT,
AWS SES, GCP, Azure DevOps, Mockito, JUnit, Karma, SonarQube, Git.}

\vspace{2pt}

%% ── 8. Zoluxiones / SURA Perú  (sep 2019 – dic 2019) ─────────
\cvEntryClient%
  {Ingeniero Full Stack}%
  {Zoluxiones --- Lima, Perú}%
  {SURA Perú}%
  {Contrato Digital}%
  {Septiembre 2019 -- Diciembre 2019}

\begin{cvItems}
  \cvItem{Sustituyó el papel por contratos digitales para corredores de inversión mediante una
    \textbf{aplicación móvil Android} construida con \textbf{Ionic} y \textbf{Angular}, más una
    plataforma web complementaria.}
  \cvItem{Desarrolló microservicios en \textbf{Java 8} con \textbf{Spring Boot} y servicios
    \textbf{serverless} en \textbf{Node.js / Express} sobre \textbf{AWS Lambda}.}
  \cvItem{Integró \textbf{AWS SES, SNS y SQS}; normalizó esquemas \textbf{MySQL} en
    \textbf{AWS RDS}.}
  \cvItem{Implementó autenticación con \textbf{Azure Active Directory} y aplicó buenas prácticas
    de seguridad para proteger la aplicación.}
  \cvItem{Construyó informes exportables con \textbf{JasperReport} y desplegó los microservicios
    con \textbf{Jenkins} en \textbf{AWS ECS}; analizó calidad con \textbf{SonarQube}.}
\end{cvItems}

{\footnotesize\color{cvMuted}\textit{Stack:}~%
Java 8, Spring Boot, Node.js, Express, Ionic, Angular 8, Angular Material, Bootstrap, MySQL,
AWS (RDS, Lambda, SES, SNS, SQS, ECS), Azure AD, Jenkins, JasperReport, Mockito, JUnit,
SonarQube, GitLab, Jira.}

\vspace{2pt}

%% ── 9. PSS Perú / BanBif  (abr – jul 2019) ───────────────────
%%     Contenido aportado por el candidato el 2026-10-26.
%%     TODO menor: fijar la versión de Java. El candidato recuerda una LTS
%%     vigente en 2019 — eran Java 8 y Java 11 — pero no cuál.
\cvEntryClient%
  {Ingeniero Full Stack}%
  {PSS Perú --- Lima, Perú}%
  {BanBif --- Banco Interamericano de Finanzas}%
  {Productos de Ahorro}%
  {Abril 2019 -- Julio 2019}

\begin{cvItems}
  \cvItem{Construyó nuevos productos de ahorro para \textbf{BanBif} --- \textbf{Ahorro Programado}
    y \textbf{Ahorro Automático} --- desde la definición del producto hasta producción.}
  \cvItem{Desarrolló servicios backend en \textbf{Java} con \textbf{Spring Boot}, exponiendo
    interfaces \textbf{REST} y \textbf{SOAP} para integrarse con los sistemas existentes del
    banco.}
  \cvItem{Mantuvo y amplió módulos web heredados construidos sobre \textbf{JSP} y \textbf{JSTL}.}
  \cvItem{Desarrolló el frontend de cara al cliente en \textbf{Angular} con
    \textbf{Angular Material}.}
\end{cvItems}

{\footnotesize\color{cvMuted}\textit{Stack:}~%
Java, Spring Boot, REST, SOAP, JSP, JSTL, Angular, Angular Material.}

\vspace{2pt}

%% ── 10. Experis / Equifax — KATRINA  (jul 2018 – abr 2019) ───
%%     Corregido. El CUL cierra el 12/04/2019; el CV anterior decía sep 2019.
\cvEntryClient%
  {Desarrollador Senior de Java}%
  {Experis / Manpower Servicios Profesionales --- Lima, Perú}%
  {Equifax}%
  {Transformación Digital (KATRINA)}%
  {Julio 2018 -- Abril 2019}

\begin{cvItems}
  \cvItem{Migró y construyó microservicios de motor de decisiones adaptados a cada cliente
    corporativo --- bancos, financieras y proveedores de crédito.}
  \cvItem{Desarrolló microservicios con \textbf{Java 8}, \textbf{Spring Boot} y
    \textbf{Spring Cloud}, exponiendo interfaces \textbf{REST} y \textbf{SOAP}.}
  \cvItem{Implementó la comunicación entre servicios con \textbf{RestTemplate} y
    \textbf{RabbitMQ}; contenerizó servicios con \textbf{Docker}.}
  \cvItem{Construyó módulos frontend en \textbf{Angular 6}, con pruebas unitarias en
    \textbf{Karma} y \textbf{Jasmine}; diseñó informes con \textbf{JasperReport}.}
  \cvItem{Entregó de forma continua con \textbf{GoCD}; aseguró la calidad del código con
    \textbf{SonarQube} y gestionó el repositorio en \textbf{SVN} y la documentación en
    \textbf{Confluence}.}
\end{cvItems}

{\footnotesize\color{cvMuted}\textit{Stack:}~%
Java 8, Spring Boot, Spring Cloud, Angular 6, Bootstrap, Docker, RabbitMQ, REST, SOAP,
JasperReport, Mockito, JUnit, Karma, SonarQube, GoCD, SVN, Confluence, Jira.}

\vspace{2pt}

%% ── 11. Olva Courier — ERP  (ago 2017 – may 2018) ────────────
%%     Corregido. El CUL registra 24/08/2017 -- 31/05/2018; el CV anterior
%%     decía mar 2015, sin respaldo de planilla. Se reescribió el bullet que
%%     hablaba de "una permanencia de tres años".
\cvEntry%
  {Desarrollador Java}%
  {Olva Courier --- Lima, Perú\enspace\textbar\enspace\textbf{Proyecto:}~ERP Olva Corp}%
  {Agosto 2017 -- Mayo 2018}

\begin{cvItems}
  \cvItem{Entregó y asumió progresivamente la propiedad de nuevos módulos del \textbf{ERP}
    corporativo de uno de los principales operadores de courier privado del Perú, abarcando
    mantenimiento, desarrollo nuevo e integración de sistemas.}
  \cvItem{Migró funcionalidades de múltiples sistemas antiguos hacia la plataforma ERP
    consolidada, retirando herramientas fragmentadas de la operación.}
  \cvItem{Desarrolló módulos del ERP en \textbf{Java 8} con \textbf{PrimeFaces}.}
  \cvItem{Construyó servicios \textbf{SOAP} y \textbf{REST} para integrar el ERP con sistemas
    internos y de socios logísticos.}
  \cvItem{Desarrolló informes operativos y financieros con \textbf{JasperReport}.}
  \cvItem{Construyó una \textbf{aplicación móvil Android} en \textbf{Kotlin} para operaciones de
    campo y reparto.}
  \cvItem{Escribió pruebas unitarias con \textbf{Mockito} y \textbf{JUnit}; desplegó en servidores
    \textbf{GlassFish} y gestionó el código en \textbf{GitLab}.}
\end{cvItems}

{\footnotesize\color{cvMuted}\textit{Stack:}~%
Java 8, PrimeFaces, SOAP, REST, JasperReport, Kotlin (Android), GlassFish, Mockito, JUnit,
GitLab, Trello.}

\vspace{2pt}

%% ── 12. Desarrollador Web Freelance  (mar 2015 – jul 2017) ───
%%
%%  Cubre la ventana entre LimaW y Olva. Decisión del candidato, 2026-10-26:
%%  Alternativa Tecnológica / BCP (dic 2016 – mar 2017) se retiró del CV y
%%  este bloque freelance cubre todo el tramo.
%%
%%  Este período NO PUEDE aparecer en el Certificado Único Laboral, y no es
%%  una omisión: las consideraciones del propio certificado establecen que no
%%  se presentan los casos de trabajadores que generan ingresos de
%%  "cuarta – quinta categoría", que es exactamente la figura del trabajo
%%  independiente. El certificado calla ahí por diseño.
%%
%%  TODO — AMPLIAR
%%  Es el bloque más largo de la carrera temprana (2 a 5 m) y hoy tiene un
%%  solo bullet. Faltan tipo y cantidad de clientes, stack concreto y algún
%%  proyecto nombrable.
\cvEntry%
  {Desarrollador Web Freelance}%
  {Independiente --- Lima, Perú}%
  {Marzo 2015 -- Julio 2017}

\begin{cvItems}
  \cvItem{Entregó desarrollo web para clientes directos como contratista independiente,
    cubriendo construcción, despliegue y mantenimiento posterior.}
\end{cvItems}

\vspace{2pt}

%% ── 13. LimaW / WProducciones — Biometría Zircon  (may 2014 – feb 2015)
%%     Entrada fusionada. Confirmado por el candidato el 2026-10-26:
%%     WProducciones y LimaW son la misma empresa. El CUL trae dos registros
%%     — WProducciones 01/05/2014 -- 31/01/2015 (9 m) y LimaW
%%     01/02/2015 -- 28/02/2015 (1 m) — que juntos dan 10 meses continuos.
\cvEntry%
  {Desarrollador Java}%
  {LimaW / WProducciones --- Lima, Perú\enspace\textbar\enspace\textbf{Proyecto:}~Biometría Zircon}%
  {Mayo 2014 -- Febrero 2015}

\begin{cvItems}
  \cvItem{Construyó una plataforma para gestionar registros de asistencia obtenidos de marcadores
    biométricos, con servicios adaptados a los sistemas de cada cliente.}
  \cvItem{Desarrolló módulos en \textbf{Java 7 / Java 8} con \textbf{JSP} y construyó servicios de
    integración \textbf{SOAP} y \textbf{REST}.}
  \cvItem{Elaboró informes específicos por cliente con \textbf{JasperReport}.}
  \cvItem{Desplegó en servidores del cliente o en instancias \textbf{JBoss} alojadas por la
    empresa según el contrato; configuró las bases \textbf{MySQL} de cada cliente y atendió
    incidencias.}
  \cvItem{Escribió pruebas unitarias con \textbf{Mockito} y \textbf{JUnit}; gestionó el código en
    \textbf{SVN} bajo Kanban.}
\end{cvItems}

{\footnotesize\color{cvMuted}\textit{Stack:}~%
Java 7/8, JSP, SOAP, REST, JasperReport, JBoss, MySQL, Mockito, JUnit, SVN.}
